\SourceLeaf{254}%
will be used here as well. For any point of the surface, \(r\) is to be its shortest distance from a fixed point, and \(\varphi\) is the angle at that fixed point between the first element of \(r\) and a fixed direction. Thus we have \(E'=1\), \(F'=0\), \(\omega'=90^0\). We also set \(\BareRoot{G'}=m\), so that any line element \mbox{is \(=\BareRoot{(\mathrm{d}r^2+mm\mathrm{d}\varphi^2)}\)}. The four equations derived in the preceding article for \(\alpha,\beta,\gamma,\delta\) therefore give:
\[
\begin{aligned}
\BareRoot{E}\mathord{.}\cos(\omega-\psi)&=\frac{\mathrm{d}r}{\mathrm{d}p}\quad .\ .\ .\ .\ .\ .\ .\ .\ .\ .\ .\ .\quad(1)\\
\BareRoot{G}\mathord{.}\cos\psi&=\frac{\mathrm{d}r}{\mathrm{d}q}\quad .\ .\ .\ .\ .\ .\ .\ .\ .\ .\ .\ .\quad(2)\\
\BareRoot{E}\mathord{.}\sin(\psi-\omega)&=m\mathord{.}\frac{\mathrm{d}\varphi}{\mathrm{d}p}\quad .\ .\ .\ .\ .\ .\ .\ .\ .\ .\ .\quad(3)\\
\BareRoot{G}\mathord{.}\sin\psi&=m\mathord{.}\frac{\mathrm{d}\varphi}{\mathrm{d}q}\quad .\ .\ .\ .\ .\ .\ .\ .\ .\ .\ .\quad(4)
\end{aligned}
\]
The last and the next-to-last equations of that article, however, give
\[
\begin{aligned}
EG-FF&=E(\frac{\mathrm{d}r}{\mathrm{d}q})^2-2F\mathord{.}\frac{\mathrm{d}r}{\mathrm{d}p}\mathord{.}\frac{\mathrm{d}r}{\mathrm{d}q}+G(\frac{\mathrm{d}r}{\mathrm{d}p})^2\quad .\ .\ .\ .\ .\ .\quad(5)\\
(E\mathord{.}\frac{\mathrm{d}r}{\mathrm{d}q}-F\mathord{.}\frac{\mathrm{d}r}{\mathrm{d}p})\mathord{.}\frac{\mathrm{d}\varphi}{\mathrm{d}q}&=(F\mathord{.}\frac{\mathrm{d}r}{\mathrm{d}q}-G\mathord{.}\frac{\mathrm{d}r}{\mathrm{d}p})\mathord{.}\frac{\mathrm{d}\varphi}{\mathrm{d}p}\quad .\ .\ .\ .\ .\ .\ .\quad(6)
\end{aligned}
\]

From these equations we must determine \(r,\varphi,\psi\) and, if it seems necessary, \(m\) in terms of \(p\) and \(q\). Integrating equation (5) gives \(r\). Once this has been found, integrating equation (6) gives \(\varphi\), and either equation (1) or (2) gives \(\psi\) itself. Finally, \(m\) is obtained from either equation (3) or (4).

The general integration of equations (5) and (6) must introduce two arbitrary functions. Their meaning is easy to understand if we consider that these equations apply just as well beyond the case considered here, when \(r\) and \(\varphi\) are taken in the more general sense of Article 16. In that case \(r\) is the length of a shortest line drawn normal to a fixed but arbitrarily chosen line, and \(\varphi\) is an arbitrary function of the length of the portion of that line between a variable shortest line and an arbitrarily chosen fixed point. The general solution must therefore embrace all these possibilities, and the arbitrary functions become determinate only once that arbitrarily chosen line and the function of its portions that \(\varphi\) is to represent have been prescribed. In our case we can take an infinitesimal circle centred at the point from which the distances \(r\) are measured, and \(\varphi\) will denote the lengths of the arcs of this circle divided by the radius,
